---
date: ``2026-07-25''
tags:
  - daily
draft: ``false''
aliases:
---

up:: Valor de Uso

\section{A satisfação de necessidades materiais tem de se dar materialmente}
A satisfação de necessidades \textit{humanas} trata-se de uma condição, em primeira instância, inorgânica/orgânica. É a necessidade de matar a sede e a fome o que mantém os animais vivos, e o ser humano não escapa desta ineliminável necessidade. 

Por vezes, se faz necessário dizer o óbvio: a satisfação de necessidades materiais tem de se dar, necessariamente, por meios materiais. Pode até ser que Siddhartha Gautama tenha alcançado o Nirvana através de meditação e da abnegação do desejo, mas não poderia fazê-lo em vida sem um cantil d'água ao seu lado; pode tê-lo feito enquanto em jejum, mas em algum momento teria de amortecer sua fome ao risco de não se tornar Buddha em vida.

!Bens imateriais também são necessariamente materiais

\subsubsection{São tanto mais valor de uso quanto menos as necessidades a se satisfazer sejam (imediatamente) [in]orgânicas}
Valores de uso, porém, como Marx o coloca logo na primeira página do livro I d'O Capital, satisfazem necessidades não só do estômago, \textit{mas também da imaginação}.\footnote{``A mercadoria é, antes de tudo, um objeto externo, uma coisa que, por meio de suas propriedades, satisfaz necessidades humanas de um tipo qualquer. A natureza dessas necessidades — se, por exemplo, elas provêm do estômago ou da imaginação — não altera em nada a questão. Tampouco se trata aqui de como a coisa satisfaz a necessidade humana, se diretamente, como meio de subsistência [\textit{Lebensmittel}], isto é, como objeto de fruição, ou indiretamente, como meio de produção.'' \parencite[p. 113]{Marx2017}} E é, em verdade, quanto mais tais necessidades não são (puramente) remetentes à esfera inorgânica/orgânica que objetos tornam-se mais e mais ``valores de uso''. Ou seja, é justamente quanto mais tais necessidades são de caráter exclusivamente social que elas serão tão mais ``valor de uso'', invés de um mero meio de satisfação de necessidades.\footnote{Em seu fascinante livro \textit{Biomimicry}, Janine Benyus comenta sobre como uma espécie de bugios (``\textit{mantled howler monkeys}'', \textit{Alouatta palliata}) conseguem induzir um viés estatístico não-desprezível no gênero de seus filhotes (``nove machos em dez, ou quatro fêmeas em cinco''). Experimentos mostraram que, a depender do que as fêmeas ingeriam, era possível induzir uma diferença de potencial elétrico mínima (na ordem de milivolts) tal que seria possível atrair ou repelir espermatozoides com cromossomos X (femininos, eletropositivo) ou Y (masculinos, eletronegativos) quando da cópula. Benyus explica por que compensaria fazê-lo: ``Se a população tem falta de machos, a fêmea que gerar machos terá uma boa chance de gerar algum que se tornará um líder do bando. Produzir um filho que seja um líder confere \textit{status} à mãe (melhor acesso a alimentação e segurança, por exemplo). Se a população tem falta de fêmeas, a mãe pode querer ter fêmeas que poderiam se tornar primeiras-damas — fazendo da mãe uma \textit{in-law of royalty}'' \parencite[p. 171]{Benyus2009}. Isso serve de exemplo de como não faltam exemplos impressionantes de ``causalidade'' na natureza, e que, não obstante isso, eles não precisem ter caráter \textit{deliberado}, ou teleológico, para que sejam assim apreciados.}

A humanização do ser humano se dá justamente quanto mais ele busca satisfazer necessidades que não sejam (imediatamente) [in]orgânicas \parencite{Lukacs2012}.
\begin{quote}
``o trabalho é antes de tudo, em termos genéticos, o ponto de partida para o tornar-se homem do homem, para a formação das suas faculdades, sendo que jamais se deve esquecer o domínio sobre si mesmo.'' \parencite[p. 348]{Lukacs2012}

``Fome é fome, mas a fome que se sacia com carne cozida, comida com garfo e faca, é uma fome diversa da fome que devora carne crua, com mão, unha e dente.'' \parencite[ ]{Marx2021}
\end{quote}

Tais necessidades
\begin{quote}
“assumem no ser social caráter completamente distinto dos processos presentes nos seres da natureza orgânica – em outras palavras, as necessidades biológicas adquirem no homem caráter eminentemente social, sem nunca abandonar em definitivo a base biológica sobre a qual sempre se apoia.” \parencite[p. 58]{Fortes2016}
\end{quote}

\section{Valores de uso são potência de satisfação de necessidades}
\begin{quote}
``O valor de uso se efetiva apenas no uso ou no consumo.'' \parencite[p. 114]{Marx2017}
\end{quote}

---
\section{O trabalho enquanto produtor de valores de uso}
\subsubsection{O trabalho enquanto condição de humanização do ser humano}
Ao longo de sua existência, o ser humano não satisfez suas necessidades somente com os objetos que encontrou prontos ao seu redor. Sua vida enquanto caçador-coletor pode ter consistido primariamente (mas não só) neste caráter de dependência dos frutos imediatos da Mãe-Natureza, mas não é à toa que o surgimento da civilização geralmente é associado com os assentamentos humanos em decorrência da agricultura: porque é então que esta espécie animal, a princípio jogada para lá e para cá pelas intempéries do mundo, passa a impor sua existência assertivamente, estabelecendo suas próprias condições de vida de maneira deliberada. É então, portanto, que ela devém algo para além de mera animalidade: é então que ela devém espécie \textit{humana}. 

\begin{quote}
``Como criador de valores de uso, como trabalho útil, o trabalho é, assim, uma condição de existência do homem, independente de todas as formas sociais, eterna necessidade natural de mediação do metabolismo entre homem e natureza'' \parencite[p. 120]{Marx2017}
\end{quote}

\subsubsection{O caráter de alternativas do trabalho}
É justamente quando passa a \textit{pôr} seus pressupostos que o ser humano devém humano. Tal pôr não consiste somente na produção, mas, e tanto mais relevantemente, na produção \textit{de meios adequados} para seus fins. O ser humano é tanto mais humano quanto mais consegue entrever, na infinitude qualitativa do mundo\footnote{\hl{Primeiro capítulo da Fenomenologia do Espírito sobre empirismo?}}, aquilo que é mais aderente a seus fins daquilo que não o é. Ou seja, é o caráter de \textit{escolha entre alternativas} aquilo que o torna mais do que mero animal no mundo. 

\subsubsection{Legalidades naturais em legalidades postas}
Objetos da natureza adquirem — ou melhor, \textbf{têm a si atribuída} — uma valoração ``de uma perspectiva externa'' \parencite[p 181]{Medeiros2016}. Objetos ``parid[o]s pelo agir humano'', na \textit{sociedade}, já surgem ``como objetivações de escolhas, como valores que se realizam'' (Ibid.). Não só objetos, mas \textit{ações} são passíveis de escolhas — logo, de valoração —; portanto, a atividade humana é, em geral, ``autoavaliada''.

Neste fazer, o ser humano altera não só a natureza, mas também \textit{sua própria} natureza: ``[atuando] sobre a natureza externa e modificando-a por meio desse movimento, ele modifica, ao mesmo tempo, sua própria natureza.'' \parencite[p. 255]{Marx2017}.

Altera não só seu caráter enquanto indivíduo humano (durante o processo de trabalho, pela mobilização de músculos etc etc), como também seu hábitat. Portanto, por meio destas condições ``pré-postas'', o trabalho atual é condicionado. (?)

---
\subsubsection{Referências}
